\documentclass[a4paper,12pt]{article}
\usepackage[top=2cm,bottom=2cm,left=2.5cm,right=1.5cm]{geometry}
\usepackage{fontspec}
\usepackage{polyglossia}
\setmainlanguage{russian}
\IfFontExistsTF{Times New Roman}{%
  \setmainfont{Times New Roman}[Ligatures=TeXOff]%
  \newfontfamily\cyrillicfont{Times New Roman}[Ligatures=TeXOff]%
}{\IfFontExistsTF{CMU Serif}{%
  \setmainfont{CMU Serif}[Ligatures=TeXOff]%
  \newfontfamily\cyrillicfont{CMU Serif}[Ligatures=TeXOff]%
}{}}
\IfFontExistsTF{Courier New}{%
  \setmonofont{Courier New}%
  \newfontfamily\cyrillicfonttt{Courier New}%
}{\IfFontExistsTF{CMU Typewriter Text}{%
  \setmonofont{CMU Typewriter Text}%
  \newfontfamily\cyrillicfonttt{CMU Typewriter Text}%
}{}}
\usepackage{amsmath,amssymb}
\usepackage{float}
\usepackage{array}
\usepackage{verbatim}
\usepackage{setspace}
\setstretch{1.1}
\setlength{\parindent}{0pt}
\setlength{\parskip}{4pt}
\renewcommand{\arraystretch}{1.25}
\setlength{\tabcolsep}{4.5pt}
\emergencystretch=3em
\newcommand{\pds}[2]{\dfrac{\partial #1}{\partial #2}}

\begin{document}

\begin{center}
{\large\bfseries Лабораторная работа\\[2pt]
Метод множителей Лагранжа}\\[4pt]
Дисциплина: Методы оптимизации
\end{center}

\begin{tabular}{@{}ll}
Ф.И.О.: & Артемов Илья Витальевич\\
Поток: & 1.1\\
Вариант задания: & 2\\
\end{tabular}

\section{Постановка задачи}

Найти условные минимум и максимум функции $f(x,y)=xy$ при ограничении $x^2+y^2=8$, где $x,y\in\mathbb R$. Это вариант 2 из раздела "Свой вариант" файла \texttt{lagrange-lecture.html}.

Требуется пройти шаги ручного решения с пары: проверить непустоту допустимого множества и регулярность ограничения, составить функцию Лагранжа с плюсом перед множителем, выписать полную систему и найти все действительные решения. Каждую точку нужно подставить в ограничение и целевую функцию, определить характер экстремума и отдельно обосновать глобальность.

Решение проверяется программой на Python с библиотекой SymPy. Программа сама дифференцирует функцию Лагранжа, составляет систему и печатает все корни, невязку ограничения и значение $f$. Ручной и программный ответы сопоставляются рядом в итоговой таблице.

\section{Ход решения}

\subsection*{Шаг 1. Функция, ограничение и допустимое множество}

Запишем цель и ограничение с нулём справа:
\[
f(x,y)=xy,\qquad h(x,y)=x^2+y^2-8=0.
\]
Допустимое множество - окружность радиуса $\sqrt8=2\sqrt2$ с центром в начале координат. Оно непусто: например, точка $(2\sqrt2;0)$ удовлетворяет ограничению, так как $(2\sqrt2)^2+0^2-8=0$.

Функции $f$ и $h$ являются многочленами и непрерывно дифференцируемы. Градиент ограничения:
\[
\nabla h=(2x;2y).
\]
Он равен нулю только при $x=y=0$, но эта точка не лежит на окружности. На всём допустимом множестве $\|\nabla h\|^2=4(x^2+y^2)=32>0$. Значит, все допустимые точки регулярны и метод Лагранжа применим.

\subsection*{Шаг 2. Число переменных и ограничений}

Исходных переменных две: $x$ и $y$, поэтому $n=2$. Ограничение одно, $m=1$. Вводим множитель $\lambda$ и получаем три неизвестных: $x$, $y$, $\lambda$. Полная система должна содержать $n+m=3$ уравнения.

\subsection*{Шаг 3. Функция Лагранжа}

Перед множителем используем знак "плюс":
\[
L(x,y,\lambda)=f+\lambda h=xy+\lambda(x^2+y^2-8).
\]
После раскрытия скобок $L=xy+\lambda x^2+\lambda y^2-8\lambda$.

\subsection*{Шаг 4. Частные производные по переменным}

При вычислении производной по $x$ переменные $y$ и $\lambda$ считаем постоянными. Из $xy$ получаем $y$, из $\lambda x^2$ получаем $2\lambda x$. Аналогично вычисляем производную по $y$:
\[
\frac{\partial L}{\partial x}=y+2\lambda x,\qquad
\frac{\partial L}{\partial y}=x+2\lambda y.
\]
Градиент целевой функции $\nabla f=(y;x)$.

\subsection*{Шаг 5. Условия стационарности}

Приравниваем производные по $x$ и $y$ к нулю:
\[
y+2\lambda x=0,\qquad x+2\lambda y=0.
\]
Это покоординатная запись условия $\nabla f+\lambda\nabla h=0$. К этим двум уравнениям нужно добавить ограничение $h=0$. Решения полной системы дают допустимые стационарные точки, а их характер проверяется отдельно.

\subsection*{Шаг 6. Ограничение и полная система}

Производная по множителю возвращает ограничение:
\[
\frac{\partial L}{\partial\lambda}=x^2+y^2-8.
\]
Получаем три уравнения для трёх неизвестных:
\[
\begin{cases}
\pds{L}{x}=y+2\lambda x=0,\\[10pt]
\pds{L}{y}=x+2\lambda y=0,\\[10pt]
\pds{L}{\lambda}=x^2+y^2-8=0.
\end{cases}
\]

\subsection*{Шаг 7. Все действительные решения системы}

Сначала проверим нулевые координаты. При $x=0$ из первого уравнения следует $y=0$, но тогда ограничение даёт $0=8$. Поэтому $x\ne0$. Аналогично при $y=0$ из второго уравнения получилось бы $x=0$, так что $y\ne0$.

Умножим первое уравнение на $y$, а второе на $x$:
\[
y^2+2\lambda xy=0,\qquad x^2+2\lambda xy=0.
\]
Вычитая второе из первого, исключаем множитель:
\[
y^2-x^2=0,\qquad (y-x)(y+x)=0.
\]
Рассмотрим обе возможности.

\textbf{Случай 1: $y=x$.} Подставим в ограничение: $2x^2=8$, откуда $x=\pm2$. Из первого уравнения $x+2\lambda x=x(1+2\lambda)=0$. Так как $x\ne0$, то $\lambda=-\frac12$. При этих $y$ и $\lambda$ второе уравнение также выполнено: $x+2\lambda y=x-x=0$. Получаем два решения $(x;y;\lambda)$:
\[
\left(-2;-2;-\frac12\right),\qquad
\left(2;2;-\frac12\right).
\]

\textbf{Случай 2: $y=-x$.} Ограничение снова даёт $2x^2=8$, $x=\pm2$. Из первого уравнения $-x+2\lambda x=x(-1+2\lambda)=0$, поэтому $\lambda=\frac12$. Второе уравнение также выполнено: $x+2\lambda y=x-x=0$. Получаем ещё два решения $(x;y;\lambda)$:
\[
\left(-2;2;\frac12\right),\qquad
\left(2;-2;\frac12\right).
\]

Все случаи рассмотрены: нулевые координаты невозможны, а равенство $y^2=x^2$ даёт только две разобранные ветви. Найдены все четыре действительных решения; при указанных множителях они удовлетворяют первым двум уравнениям системы.

\subsection*{Шаг 8. Проверка ограничения}

Подставим каждую точку в исходное ограничение:
\begin{align*}
h(-2;-2)&=(-2)^2+(-2)^2-8=4+4-8=0,\\
h(2;2)&=2^2+2^2-8=4+4-8=0,\\
h(-2;2)&=(-2)^2+2^2-8=4+4-8=0,\\
h(2;-2)&=2^2+(-2)^2-8=4+4-8=0.
\end{align*}
Все четыре точки допустимы. Невязка ограничения в каждой равна нулю.

\subsection*{Шаг 9. Характер точек и глобальность}

Допустимое множество непусто, замкнуто и ограничено, то есть компактно. Функция $f=xy$ непрерывна, поэтому по теореме Вейерштрасса на окружности существуют глобальные минимум и максимум.

В шаге 1 проверено, что все точки окружности регулярны. Поэтому в каждой точке глобального экстремума выполняются необходимые условия Лагранжа. В шаге 7 найдены все действительные решения полной системы, следовательно, все точки глобального минимума и максимума находятся среди четырёх полученных кандидатов.

Сравним значения функции во всех этих точках (подстановка приведена в шаге 10):
\[
f(-2;-2)=f(2;2)=4,\qquad f(-2;2)=f(2;-2)=-4.
\]
Наибольшее из найденных значений равно $4$, наименьшее равно $-4$. Значит,
\[
f_{\max}=4,\qquad f_{\min}=-4.
\]
Точки $(-2;-2)$ и $(2;2)$ дают глобальный максимум, точки $(-2;2)$ и $(2;-2)$ дают глобальный минимум. Таким образом, полнота набора кандидатов установлена при решении системы, существование экстремумов - по теореме Вейерштрасса, а их значения определены сравнением значений функции во всех кандидатах.

\subsection*{Шаг 10. Значения целевой функции}

Подставим координаты каждой точки именно в исходную функцию $f=xy$:
\begin{align*}
f(-2;-2)&=(-2)\cdot(-2)=4,\\
f(2;2)&=2\cdot2=4,\\
f(-2;2)&=(-2)\cdot2=-4,\\
f(2;-2)&=2\cdot(-2)=-4.
\end{align*}
Значения множителя и функции различаются: на максимуме $\lambda=-\frac12$, $f=4$; на минимуме $\lambda=\frac12$, $f=-4$.

\subsection*{Шаг 11. Смысл результата}

При фиксированной сумме квадратов $x^2+y^2=8$ крайние значения произведения достигаются при равных модулях координат, $|x|=|y|=2$. Среди этих точек одинаковые знаки координат дают максимум $4$, противоположные - минимум $-4$.

Например, в точке $(2;2)$ градиенты равны $\nabla f=(2;2)$ и $\nabla h=(4;4)$. Множитель $\lambda=-\frac12$ даёт $\lambda\nabla h=(-2;-2)$, поэтому $\nabla f+\lambda\nabla h=0$. Геометрически линия уровня $xy=4$ касается окружности в точках максимума; линия уровня $xy=-4$ касается её в точках минимума.

\subsection*{Шаг 12. Программа}

Программа \texttt{lagrange.py} проверяет решение с помощью символьного солвера SymPy. Исходные данные в коде - функции $f$, $h$ и формула $L=f+\lambda h$. Все три производные вычисляются через \texttt{sp.diff}, а полная система решается через \texttt{sp.solve}. Для каждого корня программа подставляет решение в $h$ и $f$ и печатает результат.

\begin{minipage}{\textwidth}
\textbf{Исходный код программы.}

\begingroup
\footnotesize\singlespacing
\verbatiminput{lagrange.py}
\endgroup
\end{minipage}

\textbf{Вывод программы.}

\begin{minipage}{\textwidth}
\begingroup
\footnotesize\singlespacing
\begin{verbatim}
Вариант 2
f = x*y
h = x**2 + y**2 - 8 = 0
grad h = [2*x, 2*y]
L = lam*(x**2 + y**2 - 8) + x*y
Система:
dL/dx = 2*lam*x + y = 0
dL/dy = 2*lam*y + x = 0
dL/dlam = x**2 + y**2 - 8 = 0
Число действительных решений: 4
1) x = -2, y = -2, lam = -1/2; h = 0; f = 4
2) x = 2, y = 2, lam = -1/2; h = 0; f = 4
3) x = -2, y = 2, lam = 1/2; h = 0; f = -4
4) x = 2, y = -2, lam = 1/2; h = 0; f = -4
\end{verbatim}
\endgroup
\end{minipage}

Программа нашла те же четыре решения. Во всех точках невязка $h=0$.

\newpage
\subsection*{Шаг 13. Сопоставление результатов и ответ}

В таблице~\ref{tab:comparison} ручной и программный результаты стоят рядом. В каждой тройке указаны $x$, $y$ и $\lambda$; столбец $h$ содержит невязку ограничения.

\begin{table}[H]
\centering
\caption{Ручное решение и проверка SymPy}
\label{tab:comparison}
\begin{tabular}{|c|c|c|c|c|c|c|}
\hline
 & \multicolumn{3}{c|}{Ручное решение} & \multicolumn{3}{c|}{SymPy}\\ \cline{2-7}
№ & $(x;y;\lambda)$ & $h$ & $f$ & $(x;y;\lambda)$ & $h$ & $f$\\ \hline
1 & $\left(-2;-2;-\frac12\right)$ & $0$ & $4$ & $\left(-2;-2;-\frac12\right)$ & $0$ & $4$\\ \hline
2 & $\left(2;2;-\frac12\right)$ & $0$ & $4$ & $\left(2;2;-\frac12\right)$ & $0$ & $4$\\ \hline
3 & $\left(-2;2;\frac12\right)$ & $0$ & $-4$ & $\left(-2;2;\frac12\right)$ & $0$ & $-4$\\ \hline
4 & $\left(2;-2;\frac12\right)$ & $0$ & $-4$ & $\left(2;-2;\frac12\right)$ & $0$ & $-4$\\ \hline
\end{tabular}
\end{table}

Итоговый ответ:
\[
\boxed{f_{\max}=4\text{ в точках }(-2;-2),\ (2;2),\quad \lambda=-\frac12,}
\]
\[
\boxed{f_{\min}=-4\text{ в точках }(-2;2),\ (2;-2),\quad \lambda=\frac12.}
\]
Оба экстремума глобальные. Ручной расчёт и программа дали одинаковые точки, множители, невязки и значения функции.

\section{Рефлективный вывод}

\textbf{Что оказалось полезным.} Метод позволил свести поиск экстремумов на окружности к трём уравнениям. Получились четыре стационарные точки: две дают максимум $4$, две - минимум $-4$. Программная проверка подтвердила ручной расчёт, а нулевые невязки показали, что все точки удовлетворяют ограничению.

\textbf{Трудности и как они преодолевались.}
\begin{enumerate}
\item При исключении множителя важно не потерять решения. Поэтому я сначала проверил нулевые координаты, затем рассмотрел обе ветви $y=x$ и $y=-x$ и оба знака $x$ в каждой ветви.
\item Глобальность требует отдельного обоснования. Непустота и компактность окружности вместе с непрерывностью $f$ гарантируют существование минимума и максимума. Регулярность ограничения и полное решение системы позволяют определить их сравнением значений в четырёх кандидатах.
\end{enumerate}

\end{document}
